% Generated draft: requires mathematical and source review.
\documentclass[11pt]{article}
\usepackage[margin=25mm]{geometry}
\usepackage{fontspec}
\setmainfont{Latin Modern Roman}
\usepackage{amsmath,amssymb,mathtools,unicode-math}
\setmathfont{Latin Modern Math}
\usepackage{enumitem,needspace,xurl,hyperref}
\setlength{\parindent}{0pt}
\setlength{\parskip}{5pt}
\setlength{\emergencystretch}{3em}
\newcommand{\N}{\mathbb{N}}
\newcommand{\Z}{\mathbb{Z}}
\newcommand{\Q}{\mathbb{Q}}
\newcommand{\R}{\mathbb{R}}
\newcommand{\C}{\mathbb{C}}
\newcommand{\F}{\mathbb{F}}
\newcommand{\A}{\mathbb{A}}
\newcommand{\PP}{\mathbb{P}}
\newcommand{\E}{\mathbb{E}}
\newcommand{\eps}{\varepsilon}
\newcommand{\dd}{\,\mathrm{d}}
\begin{document}
\section*{UnsolvedMath ID: 30005559}
\textbf{Source problem code:} catalogue entry
\subsection{Definitions and mathematical statement}
Let \(\mathcal M_g^{\mathrm{ct}}\) denote the moduli stack of stable curves of genus \(g\) of compact type, meaning that their dual graphs are trees. Their Jacobians are principally polarized abelian varieties, giving the Torelli morphism
\[
\operatorname{Tor}\colon \mathcal M_g^{\mathrm{ct}}\longrightarrow \mathcal A_g,
\]
where \(\mathcal A_g\) denotes the moduli stack of principally polarized abelian varieties of dimension \(g\). The source does not specify the ground field or the range of genera. Under the customary rational Chow interpretation, write
\[
R=R^*(\mathcal M_g^{\mathrm{ct}})
\subseteq A^*(\mathcal M_g^{\mathrm{ct}})\otimes\Q
\]
for the tautological ring. This coefficient and cycle-theoretic convention requires confirmation from the cited report.

For clarity, a graded pairing formulation of a Gorenstein kernel is as follows. Given a specified socle degree \(s\) and trace functional \(\varepsilon\colon R^s\to\Q\), set
\[
K^k=\{\alpha\in R^k:\varepsilon(\alpha\beta)=0
\text{ for every }\beta\in R^{s-k}\},
\qquad K=\bigoplus_{k\geq0}K^k,
\]
with \(R^j=0\) for \(j<0\). Thus \(K\) is the radical of the multiplication pairing followed by \(\varepsilon\). The supplied record does not specify the socle degree, trace, or an equivalent definition of its intended Gorenstein kernel; these data must be matched to the source before this formulation is definitive.

The source denotes elliptic-cover Noether--Lefschetz loci by \(\operatorname{NL}_d\) and asks about their Torelli pullbacks
\[
c_d=\operatorname{Tor}^*[\operatorname{NL}_d].
\]
It does not supply a definition of these loci, their ambient spaces, the permitted degrees \(d\), their cycle multiplicities, or the pullback construction. In particular, membership of every intended \(c_d\) in \(R\) must be established before treating these classes as generators in that ring.

\textbf{Question, with unresolved formulation data explicit.} Can all of the Gorenstein kernel of \(R^*(\mathcal M_g^{\mathrm{ct}})\) be explained using the classes \(c_d\)? If ``explained using'' means generation as a graded ideal, then, after specifying the preceding data, the proposed precise question is whether
\[
K=(c_d:d\in D)_R,
\]
or, equivalently, for every \(k\geq0\),
\[
K^k=\sum_{d\in D}R^{k-r_d}c_d,
\qquad r_d=\operatorname{codim}(c_d),
\]
where \(D\) is the intended set of degrees. This is a conditional formulation: the record does not distinguish ideal generation from linear spanning or another construction, and it does not determine \(D\), \(r_d\), or the intended quantifier over \(g\).

The supplied dated assessment says that particular Torelli-pullback classes are proved to lie in the kernel, while the whole kernel is not identified. It describes this question as requiring a larger leap than Questions 1 and 2. These are supplied source qualifications, not independently verified results or a resolution of the question.
\subsection{Short English statement}
The question asks whether elliptic-cover Noether--Lefschetz classes pulled back along the Torelli morphism account for the entire Gorenstein kernel of the tautological ring of compact-type curves. An ideal-generation interpretation is displayed conditionally. The precise loci, degree and genus ranges, pairing, and meaning of generation remain unspecified. The supplied status is open; this statement-only draft requires human mathematical review.
\subsection{Sources}
\begin{itemize}
\item Supplied local UnsolvedMath registry record, permanent ID \(30005559\), title ``Elliptic-Cover Cycles and the Gorenstein Kernel,'' exact source problem code \texttt{catalogue entry}, Oberwolfach Reports Open Problems collection; proposed year: 2023. No proposer is identified in the record.
\item \emph{Recent Trends in Algebraic Geometry} (2023), Oberwolfach Reports, pp.~1485--1542, DOI \texttt{10.4171/owr/2023/27}. Supplied report link: \url{https://doi.org/10.4171/owr/2023/27}.
\item Supplied literature reference: Pandharipande, ``Maps to a moving elliptic curve,'' \url{https://doi.org/10.4171/owr/2023/27}. This is the attribution attached to the reference; it does not by itself identify the proposer of Question 3.
\item Literature checked at: \texttt{2026-08-22}. This exact date records supplied dated curation, not an independently verified assessment. The supplied status is ``open.'' The assessment states that particular Torelli-pullback classes lie in the kernel but that the entire kernel has not been identified, and characterizes Question 3 as requiring a larger leap. The record requests independent verification of time-sensitive status and cited claims.
\item The registry labels the statement \texttt{corrected\_verified} and says it was reconstructed and checked against the official report. Those are supplied verification claims; no independent source consultation is asserted here.
\end{itemize}
\end{document}
